\begin{theorem}[S-43: Two semihomogeneous cohomology stages retain too many deformations]
\label{thm:semireg-two-semihomogeneous-stages}
Let $(X,\theta,F)$ be a real-quadratic abelian fourfold of type
$(2,2)$, with the notation $D,C,v_q$ of S-22 and $q>0$.
Let $G\subset X\times\widehat X$ be an effective finite subgroup
with a specified coherent translation--tensor action. Suppose a
$G$-equivariant perfect object $\widetilde E$ has exactly two
nonzero coherent cohomology sheaves, which after a common shift
are
\[
 U=\mathcal H^{-h}(E),\qquad V=\mathcal H^0(E),\qquad h\ge1.
 \tag{S43.1}
\]
Assume each of $U,V$ is a finite direct sum of simple
semihomogeneous vector bundles whose slopes $c_1(Q)/\operatorname{rk}(Q)$
belong to $F\theta$. Multiplicities are allowed, and the same
slope may occur in both sheaves. Suppose
\[
 e^B\operatorname{ch}(E)=Nv_q,
 \qquad N=\operatorname{rk}(E)\ne0.
 \tag{S43.2}
\]
The rank is allowed to change sign under the common shift in (S43.1).
Then $h$ is odd. If $s$ is the number of distinct slopes occurring
in $U$ or $V$, one has
\[
 s\ge4,\qquad
 \operatorname{ext}^1_G(E,E)\ge4s\ge16.
 \tag{S43.3}
\]
If the lift is simple and gluable in the equivariant category,
put $k=\chi_G(E,E)=4DN^2q^2/|G|\in\mathbf Z_{>0}$. Its full
degree-two semiregularity kernel then satisfies
\[
 \boxed{\dim\ker\sigma^2_{\widetilde E}
       \ge 4s-10+\left\lceil\frac{k}{2}\right\rceil
       \ge7.}
 \tag{S43.4}
\]
In particular no such two-stage object is fully semiregular.
This conclusion holds for every connecting class between the
two cohomology sheaves, including nonsplit ones, and for every
chosen finite linearization. It does not require the underlying
object to be simple or gluable.
\end{theorem}
\begin{proof}
\par\medskip\noindent\textbf{1. Odd Ext groups between different slopes}\quad
If $Q,Q'$ are simple semihomogeneous bundles, a common finite
isogeny pulls each back to a direct sum of copies of a line
bundle. The standard isogeny description is recalled at the
start of Gulbrandsen, \emph{Fourier--Mukai transforms of line
bundles on derived equivalent abelian varieties}, Section~6,
printed p.~133. The nondegenerate index theorem is Proposition~6.3,
pp.~133--134; its proof is read and imported as background.

If their slopes differ by $x\theta$ with $0\ne x\in F$, the
line bundle $L$ obtained from $Q^\vee\otimes Q'$ on this cover
has
\[
 \chi(L)=\deg(f)D\operatorname{Nm}_{F/\mathbf Q}(x)^2>0.
\]
Thus it is nondegenerate and has cohomology in just one degree.
The positive Euler characteristic makes that degree even.
Finite-etale trace splits $\mathcal O_X$ off $f_*\mathcal O_Y$,
so cohomology on $X$ is a direct summand of the pulled-back
cohomology. Therefore
\[
 \operatorname{Ext}^{2j+1}_X(Q,Q')=0
 \quad\text{whenever the slopes are different}.
 \tag{S43.5}
\]
This includes every odd degree; negative degrees and degrees
greater than four vanish for vector bundles on $X$.

\par\medskip\noindent\textbf{2. An intrinsic endpoint-map estimate}\quad
The canonical truncation triangle in the equivariant category is
\[
 U[h]\xrightarrow{i}E\xrightarrow{p}V
                  \xrightarrow{e}U[h+1].
 \tag{S43.6}
\]
For $a\in\operatorname{Ext}^1_G(U,U)$ and
$b\in\operatorname{Ext}^1_G(V,V)$ define
\[
 \delta_e(a,b)=a[h+1]e+e[1]b
       \in\operatorname{Ext}^{h+2}_G(V,U).
 \tag{S43.7}
\]
We claim
\[
 \operatorname{ext}^1_G(E,E)\ge\dim\ker\delta_e.
 \tag{S43.8}
\]
Indeed the shifted distinguished triangle may be written with
first two arrows $i[1],p[1]$ and last arrow $-e[1]$.
After rotating source and target twice, the condition
$\delta_e(a,b)=0$ is exactly the commuting-square condition
for the first two endpoint maps $b$ and $a[h+1]$.
Axiom TR3, in Stacks Definition~13.3.2, tag~0145, supplies a map
$u:E\to E[1]$ completing the morphism of triangles, so that
\[
 ui=i[1]a[h],\qquad p[1]u=bp.
 \tag{S43.9}
\]
The endpoint pair is uniquely recovered from $u$. The two
possible kernels, from the long exact Hom sequences, are
images of
\[
 \operatorname{Hom}_G(U[h],V)
 =\operatorname{Hom}_G(U[h+1],V[1])
 =\operatorname{Ext}^{-h}_G(U,V)=0.
\]
Consequently the linear space of triples $(a,u,b)$ that are
morphisms of these triangles injects into
$\operatorname{Ext}^1_G(E,E)$ and surjects onto
$\ker\delta_e$. This proves (S43.8) without invoking an
effectivity or moduli-space theorem.

\par\medskip\noindent\textbf{3. Independent scalar directions for each slope}\quad
Translations and degree-zero tensor twists preserve a bundle's
slope. The invariant-idempotent argument of S-37 therefore
gives actual equivariant decompositions
\[
 U=\bigoplus_t U_t,\qquad V=\bigoplus_t V_t
 \tag{S43.10}
\]
by their distinct slopes. A zero block is allowed in one of the
two sheaves. Each nonzero block has positive rank.
The construction uses invariant idempotent lifting in the
finite-dimensional endomorphism algebra, so it applies to
arbitrary linearizations rather than only to a chosen diagonal one.

For each slope $t$ choose an independent
$\alpha_t\in H^1(X,\mathcal O_X)$ and set
\[
 a|_{U_t}=\alpha_t\operatorname{id}_{U_t},\qquad
 b|_{V_t}=-\alpha_t\operatorname{id}_{V_t}.
 \tag{S43.11}
\]
These are invariant Ext classes, by naturality of the scalar
cohomology action. This assignment of endpoint pairs is
injective: on any nonzero block, composition with the perfect
trace is multiplication by its positive rank.

Suppose now that $h$ is odd. For different slopes $t,t'$,
the corresponding component of (S43.7) lies in the odd group
$\operatorname{Ext}^{h+2}_G(V_{t'},U_t)$, which vanishes by
(S43.5). For the same slope, the two terms cancel.
In fact $e$ has even degree $h+1$, so the scalar degree-one
cohomology action commutes with $e$ in the graded sense:
\[
 (\alpha_t\operatorname{id}_{U_t})[h+1]e_{tt}
       =e_{tt}[1](\alpha_t\operatorname{id}_{V_t}).
\]
The minus sign chosen in (S43.11) makes their sum zero.
Thus (S43.11) defines an injective $4s$-dimensional subspace
of $\ker\delta_e$, and (S43.8) gives
$\operatorname{ext}^1_G(E,E)\ge4s$.
This is why a shared slope does not create an exception: its
two endpoint blocks are moved by the same scalar direction,
with the sign required by the shifted triangle.

\par\medskip\noindent\textbf{4. The balanced character needs four slopes}\quad
The first Chern component of (S43.2) forces
$B=-c_1(E)/N\in F\theta$. Write the normalized slopes as
$x_j\theta_1+y_j\theta_2$, and use the signed ranks
of their occurrences in $V$ and $U[h]$ as weights $w_j$.
If $h$ were even, all these weights would be positive,
as would $N$. But their degree-two moment would satisfy
\[
 \sum_j w_jx_j^2=-Nq<0,
\]
which is impossible. Hence $h$ is odd.

The character identity has precisely the moment matrix
calculated in S-37:
\[
 \sum_j w_j
  \begin{pmatrix}1\\x_j\\y_j\\x_jy_j\end{pmatrix}
  \begin{pmatrix}1&x_j&y_j&x_jy_j\end{pmatrix}
     =\operatorname{diag}(N,-Nq,-Nq,Nq^2).
 \tag{S43.12}
\]
Its rank is four for either sign of $N$. Grouping equal
slopes shows that at least four distinct slopes have nonzero
net signed rank. Thus $s\ge4$, completing (S43.3).

\par\medskip\noindent\textbf{5. A quantitative full-trace obstruction}\quad
Assume now that $\widetilde E$ is simple and gluable.
Invariant CY4 Serre duality and the categorical Euler identity
of S-21 give
\[
 e_2=k-2+2e_1\ge k-2+8s,
 \qquad e_i=\operatorname{ext}^i_G(E,E).
\]
Let $b_2=\dim\ker\sigma^2_{\widetilde E}$. The full trace and
Atiyah-action maps are transpose and have common rank
$e_2-b_2$. Their composition has rank eighteen by S-22.
Hence
\[
 18\ge e_2-2b_2,
 \qquad
 b_2\ge4s-10+\left\lceil\frac{k}{2}\right\rceil.
\]
The positive Euler value $k$ is an integer, and $s\ge4$.
This proves (S43.4) and the asserted exclusion.
\end{proof}
